\begin{abstract}
Weighted aggregation of frozen large language models is often analyzed as if
each model's answer distribution were known.  A deployed ensemble instead sees
empirical answer frequencies from finitely many generations.  This creates two
independent sample sizes: \(Q\) labeled questions used to learn the ensemble
weights and \(N\) generations used to estimate each model's answer distribution
on each labeled training question.  Raw empirical risk minimization can exploit
sign errors near ties.  We study a margin-buffered estimator that counts a
training question as solved only when the empirical gold-answer margin exceeds
its finite-generation uncertainty radius.  Under a one-sided margin condition
at the best fixed weight, the clean population excess risk of the learned
weights is the sum of a standard
\(Q^{-1/2}\) generalization term and an
\(N^{-\beta/2}\) boundary-resolution term, up to logarithmic factors.  A
hard-margin two-point construction shows that a margin condition alone cannot
improve the \(Q\)-rate.  Faster \(Q\)-rates become available only after adding
identifiable risk growth and an absolute margin condition.  We give a direct
MILP implementation, preserve reproducible finite-generation diagnostics from
the earlier system, and state separately the benchmark experiment still needed
to evaluate the calibrated buffer itself.
\end{abstract}

\section{Introduction}
\label{sec:intro}

Repeated sampling has become a standard test-time tool for language models.
Self-consistency aggregates multiple reasoning paths, best-of-\(N\) procedures
select among sampled candidates, and multi-model systems combine complementary
models or agents \citep{wang2022selfconsistencyic,brown2024largelm,
jiang2023llmblenderel,wang2024mixtureofagentsel}.  These methods turn inference
compute into improved answer coverage, but they also make the statistical unit
of observation easy to blur.  An ensemble is evaluated across questions, while
the score used on any one question is itself estimated from repeated
generations.

We study this distinction for convex weighting of \(K\) frozen language
models, following the Best-of-Infinity setting of
\citet{komiyama2026best}.  Each model induces an unknown distribution over the
candidate answers to a question.  If those distributions were known, learning
the weight vector would be an ordinary supervised aggregation problem.  In
practice, the distribution of each model is replaced by vote frequencies from
\(N\) generations, and the weights are selected from \(Q\) labeled questions.
The resulting procedure has two independent sources of error: generalization
from \(Q\) questions and Monte Carlo error from \(N\) generations per model and
training question.

The second source is not merely another smooth estimation term.  Correctness is
determined by the sign of the smallest gold-versus-contender score gap.  A small
perturbation is harmless on a well-separated training question but can reverse
its empirical status near a tie.  Worse, raw empirical risk minimization
searches over weights using those noisy signs, so it can select a weight vector
that benefits from finite-\(N\) errors.  This is a structured instance of
selection under noisy estimates \citep{smith2006theoc}.

Our response is deliberately simple.  Instead of crediting every positive
empirical margin, buffered ERM credits a training question only when its margin
exceeds a uniform finite-generation radius
\(\varepsilon_N(\delta)=\sqrt{2\log(8KAQ/\delta)/N}\), where \(A\) bounds the
number of candidate answers.  Uniform concentration then places the noisy
buffered loss between two clean losses: the ordinary zero-threshold loss and a
clean loss with threshold \(2\varepsilon_N\).  This sandwich compares the
learned weights with one fixed population comparator and avoids imposing a
margin assumption at the random estimator.

Under the one-sided condition
\(\Pr\{0<M_q(w^\star)\le t\}\le C_m t^\beta\), the resulting excess-risk bound
is
\begin{equation}
R(\widehat w_\varepsilon)-R(w^\star)
\lesssim
\sqrt{\frac{(K-1)\log(eAQ)+\log(1/\delta)}{Q}}
+
C_m\left(\frac{\log(KAQ/\delta)}{N}\right)^{\beta/2}.
\label{eq:intro-rate}
\end{equation}
The first term is standard generalization over questions.  The second is the
mass of comparator-correct training questions that cannot be distinguished
from ties at resolution \(N^{-1/2}\).  It can decay faster than \(N^{-1/2}\)
because the target is a thresholded decision, not estimation of a smooth
probability.  Equation~\ref{eq:intro-rate} concerns the clean population risk
of the learned weights.  If deployment also uses finitely many fresh
generations, prediction-time sign errors form an additional term; we state the
available fixed-weight guarantee in Appendix~\ref{app:legacy-support} rather
than hiding it inside the main theorem.

The same margin exponent does \emph{not} automatically accelerate the
\(Q\)-term.  We give a two-model construction with margins of magnitude one at
each optimum, yet deciding which separated region of the simplex is better
reduces to distinguishing two nearly fair coins.  Its minimax excess risk is
of order \(Q^{-1/2}\).  This isolates the missing ingredient: a boundary
condition controls local disagreement, while a fast learning rate also needs
global identifiability.  We therefore state the faster result only under an
absolute margin condition together with local risk growth and global
separation.

Our contributions are:
\begin{enumerate}[leftmargin=*,itemsep=2pt,topsep=3pt]
\item We formulate finite-generation LLM weighting as a two-sample-size
problem and introduce a margin-buffered ERM with a direct mixed-integer linear
implementation.
\item We prove the robust bound in Eq.~(\ref{eq:intro-rate}) using only a
one-sided margin condition at the best fixed weight.
\item We prove that hard margins alone do not imply a fast \(Q\)-rate, then
derive an optional localized rate under explicit geometric identifiability.
\item We separate verified legacy diagnostics from direct evidence for the new
estimator and specify the artifact-complete experiment required to close that
gap.
\end{enumerate}

\section{Related Work}
\label{sec:related}

\paragraph{Test-time sampling and LLM ensembles.}
Self-consistency made majority aggregation of sampled reasoning paths a common
inference primitive \citep{wang2022selfconsistencyic}.  Repeated-sampling and
test-time scaling studies characterize how coverage and accuracy change with
additional inference compute \citep{brown2024largelm,schaeffer2025howdl,
snell2025scalinglt,wu2024inferencesl}.  Multi-model systems instead exploit
variation across frozen models, using ranking and fusion or layered agent
aggregation \citep{jiang2023llmblenderel,wang2024mixtureofagentsel}.
\citet{komiyama2026best} study the population, or best-of-infinity, limit of
weighted test-time ensembling.  Our question is complementary: what guarantee
survives when each population answer distribution is represented by only
\(N\) generations and the weights are fitted on only \(Q\) questions?

\paragraph{Aggregation of fixed predictors.}
Linear opinion pools and forecast combinations have a long history
\citep{stone1961opinionpool,bates1969combination,genest1986combining}.
Stacking learns how to combine fixed predictors
\citep{wolpert1992stacked,breiman1996stacked}, while weighted-majority methods
study adaptive combinations of experts \citep{littlestone1994weighted}.  Our
predictor is a linear opinion pool over answer distributions, followed by an
\(\arg\max\).  The question-dependent candidate set changes the geometry, but
on a fixed sample the relevant labelings are still induced by an arrangement
of affine hyperplanes in the weight simplex.

\paragraph{Margins, capacity, and fast rates.}
Margin distributions are central to the analysis of voting classifiers
\citep{schapire1997boostingtm,koltchinskii2001furthereo}.  VC growth bounds give
the baseline uniform deviation over questions
\citep{vapnik1971uniform,blumer1989learnability}, and Hoeffding concentration
controls the empirical answer frequencies \citep{hoeffding1963probability}.
Low-noise and margin assumptions can yield fast classification and aggregation
rates when paired with suitable localization or curvature
\citep{mammen1999smooth,tsybakov2004optimal,massart2006risk,
bartlett2005local,koltchinskii2006local,audibert2005fastlr}.  Our negative
example matters precisely here: boundary mass at an optimum does not rule out
a nearly tied, well-separated competing optimum.  The optional fast theorem
therefore derives its Bernstein condition from both absolute boundary control
and explicit risk growth, rather than attributing the faster rate to the margin
condition alone.

\paragraph{Certificates and diagnostics.}
PAC--Bayes analyses provide useful diagnostics for distributions over
predictors \citep{mcallester1999pacbayes,dziugaite2017computing}, and specialized
results connect Gibbs and majority-vote risks
\citep{viallard2021selfboundingmv,wu2021chebyshevcantellipi}.  The retained
PAC--Bayes numbers in this project are contextual diagnostics; they are not the
certificate proved here.  Our main argument instead uses a deterministic
buffered predictor, finite-sample margin concentration, and a clean
indicator-sandwich comparison.
